% ========================================
% CAPÍTULO 5: INTERPRETACIÓN Y JUICIO CRÍTICO  (criterio 5)
% ========================================

\chapter{INTERPRETACIÓN Y JUICIO CRÍTICO}
\label{ch:interpretacion}

\textbf{Qué indican los resultados.} (a) El auditor reconoce con fiabilidad los patrones para los que fue diseñado, pero el 100\,\% del banco \textbf{no} es evidencia de generalización: la prueba independiente baja el recall del AST a 0,89 y el del verificador a 0,78, y cada uno falla donde el otro acierta. La conclusión sostenible es que \textbf{la combinación es más robusta que cada parte}, no que el auditor sea casi perfecto. (b) La fuga de preprocesamiento y selección tuvo un efecto pequeño en los datasets elegidos, no porque sea inocua sino porque estos datos son «fáciles»: señal fuerte y AUC en el techo. El control positivo muestra que, cuando la señal es débil frente al ruido, el mismo mecanismo infla la métrica hasta 0,20--0,26; la gravedad depende de la relación entre señal y ruido y no es una propiedad fija de cada tipo de fuga ni de cada familia de modelos. (c) De las dos hipótesis de la reformulación, la del escalado en árboles se confirma con claridad; la de que la red neuronal es la más vulnerable tiene respaldo débil: se observa en regresión (\DOpt{} medio 0,017; $d=0{,}65$) pero no en clasificación, y no resiste la corrección por multiplicidad.

\textbf{Implicaciones para la factibilidad.} El proyecto es factible como instrumento (detectar, reparar y medir) y se puede defender con lo hecho. Es menos sólido como afirmación empírica sobre «qué modelos sufren más», porque hay solo 3 datasets por tarea y 9 pares por familia. Por eso el objetivo OE3 se reformula para Capstone~II: estimar \DOpt{} \textbf{condicionado a la razón señal/ruido}, y no como un número único por familia.

\textbf{Limitaciones y amenazas a la validez.} \emph{Datos:} seis datasets clásicos que los LLM probablemente conocen, un submuestreo de \texttt{california\_housing} y un efecto techo en clasificación. \emph{Método:} el generador y el auditor tienen el mismo autor (circularidad), la prueba independiente tiene 13 casos, la exclusión de un par degenerado fue posterior a ver el problema y los valores $p$ no están corregidos. \emph{Generalización:} los resultados valen para pipelines sintéticos de scikit-learn; \textbf{aún no se sabe qué tan frecuente es la fuga en código escrito por LLM reales}, que es la pregunta motivadora del proyecto, ni cómo se comporta el auditor ante su variedad de estilos.

\textbf{Riesgos para Capstone~II y ajustes propuestos.} La Tabla~\ref{tab:riesgos} los resume junto con la evidencia que los respalda.

\begin{table}[!htbp]
\centering
\caption{Riesgos observados, evidencia y ajustes para Capstone~II.}
\label{tab:riesgos}
\begingroup\footnotesize\singlespacing
\begin{tabular}{@{}p{4.4cm} p{4.0cm} p{6.6cm}@{}}
\toprule
Riesgo & Evidencia & Ajuste \\
\midrule
El código de LLM reales es más variado que las plantillas y el recall caerá & AST 0,89 y ejecución 0,78 en código escrito a mano & Ampliar la prueba independiente a $\geq 60$ scripts y, con los LLM reales, etiquetar una muestra a mano con revisión cruzada \\
Alias, funciones auxiliares y reparaciones manuales (solo 4/9 reparados) & Un caso con alias; cuatro casos sin reparación automática & Resolver alias de importación; ampliar el refactorizador a \texttt{ColumnTransformer} y funciones \\
Efecto techo y pocos datasets & AUC mediano 0,991; $|\DOpt|\leq 0{,}0034$ en clasificación & Incorporar datasets de señal débil o con muchas más variables que filas; informar \DOpt{} por tramo de señal/ruido \\
Ejecutar código de LLM es riesgoso & --- & Ejecutar solo en contenedor sin red ni credenciales; la plataforma web seguirá sin ejecutar código \\
Costo y términos de servicio de las APIs & Script de generación sin probar & Medir tokens con 2 modelos antes de escalar; revisar términos de publicación de salidas \\
\bottomrule
\end{tabular}
\endgroup
\end{table}
